\subsection{DECOMPOSITION OF A FAMILY OF ENDOMORPHISMS}

Let V be a vector space, S a set, and r a map from S to End(V). Denote by P the set of maps from S to k. If \(\lambda \in P\) , denote by \(\mathrm{V}_{\lambda}(\mathrm{S})\) (resp. \(\mathrm{V}^{\lambda}(\mathrm{S})\) ) the set of \(v \in V\) such that, for all \(s \in S\) , \(r(s)v = \lambda(s)v\) (resp. \((r(s) - \lambda(s))^{n}v = 0\) for n sufficiently large). The sets \(\mathrm{V}_{\lambda}(\mathrm{S})\) and \(\mathrm{V}^{\lambda}(\mathrm{S})\) are vector subspaces of V, and \(\mathrm{V}_{\lambda}(\mathrm{S}) \subset \mathrm{V}^{\lambda}(\mathrm{S})\) . We say that \(\mathrm{V}_{\lambda}(\mathrm{S})\) is the eigenspace of V relative to \(\lambda\) (and to r), that \(\mathrm{V}^{\lambda}(\mathrm{S})\) is the primary subspace of V relative to \(\lambda\) (and to r), and that \(\mathrm{V}^{0}(\mathrm{S})\) is the nilspace of V (relative to r). We say that \(\lambda\) is a weight of S in V if \(\mathrm{V}^{\lambda}(\mathrm{S}) \neq 0\) .

In particular, if S reduces to a single element s, P can be identified with k; we use the notations \(\mathrm{V}_{\lambda(s)}(s)\) and \(\mathrm{V}^{\lambda(s)}(s)\) , or \(\mathrm{V}_{\lambda(s)}(r(s))\) and \(\mathrm{V}^{\lambda(s)}(r(s))\) , instead of \(\mathrm{V}_{\lambda}(\{s\})\) , \(\mathrm{V}^{\lambda}(\{s\})\) ; we speak of eigenspaces, primary subspaces and the nilspace of \(r(s)\) ; an element v of \(\mathrm{V}_{\lambda(s)}(s)\) is called an eigenvector of \(r(s)\) , and, if \(v \neq 0\) , \(\lambda(s)\) is called the corresponding eigenvalue (cf.~Algebra, Chap. VII, §5).

For all \(\lambda\in P\) , the following relations are immediate:

\[
V ^ {\lambda} (S) = \bigcap_ {s \in S} V ^ {\lambda (s)} (s),\tag{1}
\]

\[
V _ {\lambda} (S) = \bigcap_ {s \in S} V _ {\lambda (s)} (s).\tag{2}
\]

Let \(k'\) be an extension of \(k\) . The canonical map from \(\operatorname{End}(V)\) to \(\operatorname{End}(V \otimes_k k')\) gives, by composition with \(r\) , a map \(r' : S \to \operatorname{End}(V \otimes_k k')\) . Similarly, every map \(\lambda\) from S to \(k\) defines canonically a map, also denoted by \(\lambda\) , from S to \(k'\) . With these notations, we have the following proposition:

PROPOSITION 1. For all \(\lambda \in \mathrm{P}\) ,

\[
(\mathrm{V} \otimes_ {k} k ^ {\prime}) ^ {\lambda} (\mathrm{S}) = \mathrm{V} ^ {\lambda} (\mathrm{S}) \otimes_ {k} k ^ {\prime} \text {and} (\mathrm{V} \otimes_ {k} k ^ {\prime}) _ {\lambda} (\mathrm{S}) = \mathrm{V} _ {\lambda} (\mathrm{S}) \otimes_ {k} k ^ {\prime}.
\]

Let \((a_{i})\) be a basis of the k-vector space \(k'\) . If \(v \in V \otimes_{k} k'\) , v can be expressed uniquely in the form \(\sum v_{i} \otimes a_{i}\) where \((v_{i})\) is a finitely-supported family of elements of V. For all \(s \in S\) ,

\[
(r ^ {\prime} (s) - \lambda (s)) ^ {n} (v) = \sum (r (s) - \lambda (s)) ^ {n} v _ {i} \otimes a _ {i}.
\]

It follows that

\[
v \in (\mathrm{V} \otimes_ {k} k ^ {\prime}) ^ {\lambda} (\mathrm{S}) \Longleftrightarrow v _ {i} \in \mathrm{V} ^ {\lambda} (\mathrm{S}) \text { for   all } i,
\]

\[
v \in (\mathrm{V} \otimes_ {k} k ^ {\prime}) _ {\lambda} (\mathrm{S}) \Longleftrightarrow v _ {i} \in \mathrm{V} _ {\lambda} (\mathrm{S}) \text { for   all } i,
\]

which implies the proposition.

PROPOSITION 2. Let \(\mathrm{V}, \mathrm{V}'\) , \(\mathrm{W}\) be vector spaces. Let \(r: \mathrm{S} \to \operatorname{End}(\mathrm{V})\) , \(r': \mathrm{S} \to \operatorname{End}(\mathrm{V}')\) and \(q: \mathrm{S} \to \operatorname{End}(\mathrm{W})\) be maps.

\begin{enumerate}
\def\labelenumi{(\roman{enumi})}
\item
  Let \(f: \mathrm{V} \to \mathrm{W}\) be a linear map such that \(q(s)f(v) = f(r(s)v)\) for \(s \in \mathrm{S}\) and \(v \in \mathrm{V}\) . Then, for all \(\lambda \in \mathrm{P}\) , \(f\) maps \(\mathrm{V}^{\lambda}(\mathrm{S})\) (resp. \(\mathrm{V}_{\lambda}(\mathrm{S})\) ) into \(W^{\lambda}(\mathrm{S})\) (resp. \(W_{\lambda}(\mathrm{S})\) ).
\item
  Let \(\mathrm{B}:\mathrm{V}\times \mathrm{V}'\to \mathrm{W}\) be a bilinear map such that
\end{enumerate}

\[
q (s) \mathrm{B} (v, v ^ {\prime}) = \mathrm{B} (r (s) v, v ^ {\prime}) + \mathrm{B} (v, r ^ {\prime} (s) v ^ {\prime})
\]

for \(s \in \mathrm{S}\) , \(v \in \mathrm{V}\) , \(v' \in \mathrm{V}'\) . Then, for all \(\lambda, \mu \in \mathrm{P}\) , B maps \(V^{\lambda}(\mathrm{S}) \times \mathrm{V}'^{\mu}(\mathrm{S})\) (resp. \(\mathrm{V}_{\lambda}(\mathrm{S}) \times \mathrm{V}_{\mu}'(\mathrm{S})\) ) into \(W^{\lambda + \mu}(\mathrm{S})\) (resp. \(W_{\lambda + \mu}(\mathrm{S})\) ).

\begin{enumerate}
\def\labelenumi{(\roman{enumi})}
\setcounter{enumi}{2}
\tightlist
\item
  Let \(B: V \times V' \to W\) be a bilinear map such that
\end{enumerate}

\[
q (s) \mathrm{B} (v, v ^ {\prime}) = \mathrm{B} (r (s) v, r ^ {\prime} (s) v ^ {\prime})
\]

for \(s \in S\) , \(v \in V\) , \(v' \in V'\) . Then, for all \(\lambda, \mu \in P\) , B maps \(V^{\lambda}(S) \times V'^{\mu}(S)\) (resp. \(V_{\lambda}(S) \times V'_{\mu}(S)\) ) into \(W^{\lambda\mu}(S)\) (resp. \(W_{\lambda\mu}(S)\) ).

In case (i), \((q(s) - \lambda(s))^n f(v) = f((r(s) - \lambda(s))^n v)\) for \(s \in S\) and \(v \in V\) , hence the conclusion. In case (ii),

\[
(q (s) - \lambda (s) - \mu (s)) \mathrm{B} (v, v ^ {\prime}) = \mathrm{B} ((r (s) - \lambda (s)) v, v ^ {\prime}) + \mathrm{B} (v, (r ^ {\prime} (s) - \mu (s)) v ^ {\prime})
\]

for \(s \in S\) , \(v \in V\) , \(v' \in V'\) , hence by induction on n

\[
(q (s) - \lambda (s) - \mu (s)) ^ {n} \mathrm{B} (v, v ^ {\prime}) = \sum_ {i + j = n} \binom {n} {i} \mathrm{B} ((r (s) - \lambda (s)) ^ {i} v, (r ^ {\prime} (s) - \mu (s)) ^ {j} v ^ {\prime}).
\]

The assertions in (ii) follow immediately. In case (iii),

\[
(q (s) - \lambda (s) \mu (s)) \mathrm{B} (v, v ^ {\prime}) = \mathrm{B} ((r (s) - \lambda (s)) v, r ^ {\prime} (s) v ^ {\prime}) + \mathrm{B} (\lambda (s) v, (r ^ {\prime} (s) - \mu (s)) v ^ {\prime})
\]

for \(s \in S\) , \(v \in V\) , \(v' \in V'\) , hence by induction on n

\[
\begin{array}{l} (q (s) - \lambda (s) \mu (s)) ^ {n} \mathrm{B} (v, v ^ {\prime}) \\ = \sum_ {i + j = n} \binom {n} {i} \mathrm{B} (\lambda (s) ^ {j} (r (s) - \lambda (s)) ^ {i} v, r ^ {\prime} (s) ^ {i} (r ^ {\prime} (s) - \mu (s)) ^ {j} v ^ {\prime}). \end{array}
\]

The assertions in (iii) follow immediately.

PROPOSITION 3. The sum \(\sum_{\lambda \in \mathrm{P}}\mathrm{V}^{\lambda}(\mathrm{S})\) is direct. The sum \(\sum_{\lambda \in \mathrm{P}}\mathrm{V}_{\lambda}(\mathrm{S})\) is direct.

The second assertion is a consequence of the first; hence it suffices to prove that. We distinguish several cases.

\begin{enumerate}
\def\labelenumi{\alph{enumi})}
\item
  S is empty. The assertion is trivial.
\item
  S is reduced to a single element s. Let \(\lambda_{0}, \lambda_{1}, \ldots, \lambda_{n}\) be distinct elements of k. For \(i = 0, 1, \ldots, n\) , let \(v_{i} \in V^{\lambda_{i}}(s)\) and assume that \(v_{0} = v_{1} + \cdots + v_{n}\) . It suffices to prove that \(v_{0} = 0\) . For \(i = 0, \ldots, n\) , there exists an integer \(q_{i} > 0\) such that \((r(s) - \lambda_{i})^{q_{i}} v_{i} = 0\) . Consider the polynomials \(\mathrm{P}(\mathrm{X}) = \prod_{i \geq 1} (\mathrm{X} - \lambda_{i})^{q_{i}}\) and \(\mathrm{Q}(\mathrm{X}) = (\mathrm{X} - \lambda_{0})^{q_{0}}\) . We have \(\mathrm{Q}(r(s)) v_{0} = 0\) , and \(\mathrm{P}(r(s)) v_{0} = \sum_{i=1}^{n} \mathrm{P}(r(s)) v_{i} = 0\) . Since P and Q are relatively prime, the Bezout identity proves that \(v_{i} = 0\) .
\end{enumerate}

\[
v _ {0} = 0
\]

\begin{enumerate}
\def\labelenumi{\alph{enumi})}
\setcounter{enumi}{2}
\item
  S is finite and non-empty. We argue by induction on the cardinal of S. Let \(s \in \mathrm{S}\) and \(\mathrm{S}' = \mathrm{S} - \{s\}\) . Let \((v_{\lambda})_{\lambda \in \mathrm{P}}\) be a finitely-supported family of elements of V such that \(\sum_{\lambda \in \mathrm{P}} v_{\lambda} = 0\) and \(v_{\lambda} \in \mathrm{V}^{\lambda}(\mathrm{S})\) . Let \(\lambda_0 \in \mathrm{P}\) . Let \(\mathrm{P}'\) be the set of \(\lambda \in \mathrm{P}\) such that \(\lambda |\mathrm{S}' = \lambda_0|\mathrm{S}'\) . By the induction hypothesis applied to \(\mathrm{S}'\) , we have \(\sum_{\lambda \in \mathrm{P}'} v_{\lambda} = 0\) . If \(\lambda, \mu\) are distinct elements of \(\mathrm{P}'\) , \(\lambda(s) \neq \mu(s)\) . Since the sum \(\sum_{\alpha \in k} \mathrm{V}^{\alpha}(s)\) is direct by \(b\) ), and since \(v_{\lambda} \in \mathrm{V}^{\lambda(s)}(s)\) , \(v_{\lambda} = 0\) for all \(\lambda \in \mathrm{P}'\) , and in particular \(v_{\lambda_0} = 0\) , which we had to prove.
\item
  General case. Let \((v_{\lambda})_{\lambda\in\mathrm{P}}\) be a finitely-supported family of elements of V such that \(\sum_{\lambda\in\mathrm{P}} v_{\lambda}=0\) and \(v_{\lambda}\in\mathrm{V}^{\lambda}(\mathrm{S})\) . Let \(P'\) be the finite set of \(\lambda\in P\) such that \(v_{\lambda}\neq0\) , and let \(S'\) be a finite subset of S such that the conditions \(\lambda\in P'\) , \(\mu\in P'\) , \(\lambda|S'= \mu|S'\) imply that \(\lambda=\mu\) . We have \(v_{\lambda}\in\mathrm{V}^{\lambda|S'}(S')\) ; applying c), we see that \(v_{\lambda}=0\) for \(\lambda\in P'\) , which completes the proof.
\end{enumerate}

Recall that, if \(x \in \operatorname{End}(\mathrm{V})\) , we denote by ad \(x\) the map \(y \mapsto xy - yx = [x, y]\) from \(\operatorname{End}(\mathrm{V})\) to itself.

Lemma 1. Let \(x, y \in \operatorname{End}(\mathbf{V})\) .

\begin{enumerate}
\def\labelenumi{(\roman{enumi})}
\item
  Assume that \(\mathrm{V}\) is finite dimensional. Then \(x\) is triangularizable if and only if \(\mathrm{V} = \sum_{a\in k}\mathrm{V}^a (x)\) .
\item
  If there exists an integer \(n\) such that \((\mathrm{ad}x)^n y = 0\) , each \(\mathrm{V}^a(x)\) is stable under \(y\) .
\item
  Assume that \(\mathrm{V}\) is finite dimensional. If \(\mathrm{V} = \sum_{a \in k} \mathrm{V}^a(x)\) and if each \(\mathrm{V}^a(x)\) is stable under \(y\) , there exists an integer \(n\) such that \((\operatorname{ad} x)^n y = 0\) .
\end{enumerate}

Part (i) follows from Algebra, Chap. VII, §5, no. 2, Prop. 3.

Let \(\mathrm{E} = \operatorname{End}(\mathrm{V})\) . Let B be the bilinear map \((u, v) \mapsto u(v)\) from \(\mathrm{E} \times \mathrm{V}\) to V. By the definition of ad \(x\) ,

\[
x (\mathrm{B} (u, v)) = B (u, x (v)) + \mathrm{B} ((\operatorname{ad} x) (u), v)
\]

for \(x \in \mathrm{E}\) , \(u \in \mathrm{E}\) , \(v \in \mathrm{V}\) . Let \(x\) operate on \(\mathrm{E}\) via \(\operatorname{ad} x\) . By Prop. 2 (ii), \(\mathrm{B}(\mathrm{E}^0(x), \mathrm{V}^a(x)) \subset \mathrm{V}^a(x)\) for all \(a \in k\) . If \((\operatorname{ad} x)^n y = 0\) , then \(y \in \mathrm{E}^0(x)\) , so \(y(\mathrm{V}^a(x)) \subset \mathrm{V}^a(x)\) , which proves (ii).

To prove (iii), we can replace V by \(V^{a}(x)\) and x (resp. y) by its restriction to \(V^{a}(x)\) . Replacing x by x - a, we can assume that x is nilpotent. Then, \((\operatorname{ad} x)^{2\dim V-1} = 0\) (Chap. I, §4, no. 2), which proves (iii).

Remark. The argument proves that, if V is finite dimensional and if there exists an integer n such that \((\operatorname{ad} x)^{n}y = 0\) , then \((\operatorname{ad} x)^{2\dim V-1}y = 0\) .

In the sequel, we shall say that the map \(r : S \to \text{End}(V)\) satisfies condition (AC) (``almost commutative'') if:

(AC) For every pair \((s, s')\) of elements of S, there exists an integer \(n\) such that

\[
(\operatorname{ad} r (s)) ^ {n} r \left(s ^ {\prime}\right) = 0.
\]

THEOREM 1. Assume that V is finite dimensional. The following conditions are equivalent:

\begin{enumerate}
\def\labelenumi{(\roman{enumi})}
\item
  Condition (AC) is satisfied and, for all \(s \in \mathrm{S}\) , \(r(s)\) is triangularizable.
\item
  For all \(\lambda \in \mathrm{P}\) , \(\mathrm{V}^{\lambda}(\mathrm{S})\) is stable under \(r(\mathrm{S})\) , and \(\mathrm{V} = \sum_{\lambda \in \mathrm{P}} \mathrm{V}^{\lambda}(\mathrm{S})\) .
\end{enumerate}

If \(V = \sum_{\lambda \in P} V^{\lambda}(S)\) , then \(V = \sum_{a \in k} V^{a}(s)\) for all \(s \in S\) , and it follows from Lemma 1 that (ii) implies (i). Assume that condition (i) is satisfied. Lemma 1 and formula (1) imply that each \(V^{\lambda}(S)\) is stable under \(r(S)\) . It remains to prove that \(V = \sum_{\lambda \in P} V^{\lambda}(S)\) . We argue by induction on dim V. We distinguish two cases.

\(a)\) For all \(s \in \mathrm{S}\) , \(r(s)\) has a single eigenvalue \(\lambda(s)\) . Then \(\mathrm{V} = \mathrm{V}^{\lambda}(\mathrm{S})\) .

\begin{enumerate}
\def\labelenumi{\alph{enumi})}
\setcounter{enumi}{1}
\tightlist
\item
  There exists \(s \in S\) such that \(r(s)\) has at least two distinct eigenvalues. Then V is the direct sum of the \(\mathrm{V}^{a}(s)\) for \(a \in k\) , and \(\dim \mathrm{V}^{a}(s) < \dim \mathrm{V}\) for all a. Each \(\mathrm{V}^{a}(s)\) is stable under \(r(\mathrm{S})\) , and it suffices to apply the induction hypothesis.
\end{enumerate}

COROLLARY 1. Assume that V is finite dimensional and that condition (AC) is satisfied. Let \(k'\) be an extension of \(k\) . Assume that, for all \(s \in S\) , the endomorphism \(r(s) \otimes 1\) of \(V \otimes_k k'\) is triangularizable. Let \(P'\) be the set of maps from \(S\) to \(k'\) . Then \(V \otimes_k k' = \sum_{\lambda' \in P'} (V \otimes_k k')^{\lambda'}(S)\) .

Let \(r': S \to \operatorname{End}(V \otimes_k k')\) be the map defined by \(r\) . If \(s_1, s_2 \in S\) , there exists an integer \(n\) such that \((\operatorname{ad} r(s_1))^n r(s_2) = 0\) , hence \((\operatorname{ad} r'(s_1))^n r'(s_2) = 0\) . It now suffices to apply Th. 1.

COROLLARY 2. Assume that V is finite dimensional and that condition (AC) is satisfied. Denote by \(\mathrm{V}^{+}(\mathrm{S})\) the vector subspace \(\sum_{s\in \mathrm{S}}\left(\bigcap_{i\geq 1}r(s)^{i}\mathrm{V}\right)\) . Then:

\begin{enumerate}
\def\labelenumi{(\roman{enumi})}
\item
  \(\mathrm{V}^0 (\mathrm{S})\) and \(\mathrm{V^{+}(S)}\) are stable under \(r(\mathrm{S})\) ;
\item
  \(\mathrm{V} = \mathrm{V}^{0}(\mathrm{S}) \oplus \mathrm{V}^{+}(\mathrm{S});\)
\item
  every vector subspace W of V, stable under \(r(\mathrm{S})\) and such that \(\mathrm{W}^0 (\mathrm{S}) = 0\) , is contained in \(\mathrm{V}^{+}(\mathrm{S})\) ;
\item
  \(\sum_{s\in S}r(s)\mathrm{V}^{+}(\mathrm{S}) = \mathrm{V}^{+}(\mathrm{S}).\)
\end{enumerate}

Moreover, \(V^{+}(S)\) is the only vector subspace of V with properties (i) and (ii). For any extension \(k'\) of k, \((\mathrm{V} \otimes_{k} k')^{+}(\mathrm{S}) = \mathrm{V}^{+}(\mathrm{S}) \otimes_{k} k'\) .

The last assertion is immediate. Thus, taking Prop. 1 into account, in proving the others we can assume that k is algebraically closed. By Th. 1, \(V = \sum_{\lambda \in P} V^{\lambda}(S)\) , and the \(V^{\lambda}(S)\) are stable under \(r(S)\) . If \(s \in S\) , the characteristic polynomial of \(r(s)|V^{\lambda}(S)\) is \((X - \lambda(s))^{\dim V^{\lambda}(S)}\) ; it follows that \(\bigcap_{i \geq 1} r(s)^{i} V^{\lambda}(s)\) is zero if \(\lambda(s) = 0\) and is equal to \(V^{\lambda}(S)\) if \(\lambda(s) \neq 0\) ; hence,

\[
\mathrm{V} ^ {+} (\mathrm{S}) = \sum_ {\lambda \in \mathrm{P}, \lambda \neq 0} \mathrm{V} ^ {\lambda} (\mathrm{S}),\tag{3}
\]

which proves (i), (ii) and (iv). If W is a vector subspace of V stable under \(r(S)\) , then \(W = \sum_{\lambda \in P} W^{\lambda}(S)\) and \(W^{\lambda}(S) = W \cap V^{\lambda}(S)\) . If \(W^{0}(S) = 0\) , we see that \(W \subset V^{+}(S)\) , which proves (iii).

Let \(V'\) be a vector subspace of \(V\) stable under \(r(S)\) and such that \(V' \cap V^0(S) = 0\) . Then \(V'^0(S) = 0\) , so \(V' \subset V^+(S)\) by (iii). If, in addition, \(V = V^0(S) + V'\) , we see that \(V' = V^+(S)\) . Q.E.D.

We sometimes call \((\mathrm{V}^{0}(\mathrm{S}),\mathrm{V}^{+}(\mathrm{S}))\) the Fitting decomposition of V, or of the map \(r:S\to\operatorname{End}(\mathrm{V})\) . If S reduces to a single element s, we write \(\mathrm{V}^{+}(s)\) or \(\mathrm{V}^{+}(r(s))\) instead of \(\mathrm{V}^{+}(\{s\})\) . We have that \(\mathrm{V}=\mathrm{V}^{0}(s)\oplus\mathrm{V}^{+}(s)\) , \(\mathrm{V}^{0}(s)\) and \(\mathrm{V}^{+}(s)\) are stable under \(r(s)\) , \(r(s)|\mathrm{V}^{0}(s)\) is nilpotent and \(r(s)|\mathrm{V}^{+}(s)\) is bijective.

COROLLARY 3. Let V and V' be finite dimensional vector spaces, and let \(r: S \to \text{End}(V)\) and \(r': S \to \text{End}(V')\) be maps satisfying condition (AC). Let \(f: V \to V'\) be a surjective linear map such that \(f(r(s)v) = r'(s)f(v)\) for \(s \in S\) and \(v \in V\) . Then \(f(V^{\lambda}(S)) = V'^{\lambda}(S)\) for all \(\lambda \in P\) .

In view of Prop. 1, we are reduced to the case in which \(k\) is algebraically closed. We have \(\mathrm{V} = \bigoplus_{\lambda \in \mathrm{P}} \mathrm{V}^{\lambda}(\mathrm{S})\) , \(\mathrm{V}' = \bigoplus_{\lambda \in \mathrm{P}} \mathrm{V}'^{\lambda}(\mathrm{S})\) by Th. 1, and \(\mathrm{V}' = f(\mathrm{V}) = \sum_{\lambda \in \mathrm{P}} f(\mathrm{V}^{\lambda}(\mathrm{S}))\) . Finally, \(f(\mathrm{V}^{\lambda}(\mathrm{S})) \subset \mathrm{V}'^{\lambda}(\mathrm{S})\) by Prop. 2 (i), hence the corollary.

PROPOSITION 4. Assume that \(k\) is perfect. Let \(V\) be a finite dimensional vector space, \(u\) an element of \(\operatorname{End}(V)\) , \(u_s, u_n\) the semi-simple and nilpotent components of \(u\) (Algebra, Chap. VII, §5, no. 8).

\begin{enumerate}
\def\labelenumi{(\roman{enumi})}
\item
  For all \(\lambda \in k\) , \(V^{\lambda}(u) = V^{\lambda}(u_s) = V_{\lambda}(u_s)\) .
\item
  If \(\mathrm{V}\) has an algebra structure and if \(u\) is a derivation of \(\mathrm{V}\) , \(u_{s}\) and \(u_{n}\) are derivations of \(\mathrm{V}\) .
\item
  If \(\mathrm{V}\) has an algebra structure and if \(u\) is an automorphism of \(\mathrm{V}\) , then \(u_{s}\) and \(1 + u_{s}^{-1}u_{n}\) are automorphisms of \(\mathrm{V}\) .
\end{enumerate}

In view of Prop. 1, we can assume that \(k\) is algebraically closed, so

\[
\mathrm{V} = \sum_ {\lambda \in k} \mathrm{V} ^ {\lambda} (u).
\]

The semi-simple component of \(u|V^{\lambda}(u)\) is the homothety with ratio \(\lambda\) in \(V^{\lambda}(u)\) . This proves (i).

Assume from now on that \(\mathrm{V}\) has an algebra structure. Let \(x \in \mathrm{V}^{\lambda}(u)\) , \(y \in \mathrm{V}^{\mu}(u)\) .

If \(u\) is a derivation of \(\mathrm{V}\) , then \(xy \in \mathrm{V}^{\lambda + \mu}(u)\) (Prop. 2 (ii)), so

\[
u _ {s} (x y) = (\lambda + \mu) (x y) = (\lambda x) y + x (\mu y) = (u _ {s} x) y + x (u _ {s} y).
\]

This proves that \(u_{s}\) is a derivation of V. Then \(u_{n} = u - u_{s}\) is a derivation of V.

If u is an automorphism of V, \(\operatorname{Ker}(u_{s}) = \mathrm{V}^{0}(u) = 0\) , so \(u_{s}\) is bijective. On the other hand, \(xy \in \mathrm{V}^{\lambda\mu}(u)\) (Prop. 2 (iii)), so

\[
u _ {s} (x y) = (\lambda \mu) (x y) = (\lambda x) (\mu y) = u _ {s} (x). u _ {s} (y).
\]

This proves that \(u_{s}\) is an automorphism of V; but then so is

\[
1 + u _ {s} ^ {- 1} u _ {n} = u _ {s} ^ {- 1} u.
\]

